% Problema 4, parte 6: campi da incollare nella piattaforma (uno per sezione)

% ===== CAMPO: 1. Result and scope =====
A complete, self-contained proof (using only the CRT criterion (*)) that
F(g) --- the maximal size of a pairwise disjoint family with all
pairwise gcds \textless= g --- satisfies the prime-splitting recursion
of Lemma 4, and the closed bounds F(g) \textless= g\^{}3 and F(g)
\textless= g\^{}2 prod\_\{p\textless=g\}(1+2/p). Consequences: for every
pairwise disjoint family of size k, max gcd \textgreater= k\^{}\{1/3\}
unconditionally, and max gcd \textgreater= c sqrt(k)/log k with Mertens'
theorem (cited). The recursion also re-proves cell 1 (R(2,2)=2). This is
strictly weaker than the cited state of the art (k\^{}\{1-o(1)\},
arXiv:2607.24655) and does NOT settle any of (a),(b),(c); it is offered
as a verified, reusable lemma and as a precise diagnosis of where the
loss occurs (the union bound over primes in Lemma 4, whose exact values
are \textasciitilde g\^{}2/2). No cell is claimed solved.


% ===== CAMPO: 2. Proof =====
\subparagraph{Setting and notation}\label{setting-and-notation}

A \emph{family} is a finite list \(\mathcal F=\{(a_i,m_i)\}_{i\in I}\)
of congruence classes \(a_i \pmod{m_i}\), \(m_i\ge 1\). It is
\emph{pairwise disjoint} if
\(a_i \pmod{m_i}\cap a_j\pmod{m_j}=\emptyset\) for all \(i\ne j\). We
use only the problem's fact \((*)\): two classes meet iff
\(\gcd(m_i,m_j)\mid a_i-a_j\).

For an integer \(g\ge 1\) and a prime \(P\) define
\[S(g,P)=\sup\Big\{|\mathcal F| :\ \mathcal F \text{ pairwise disjoint},\ \gcd(m_i,m_j)\le g\ \forall i\ne j,\ \gcd(m_i,q)=1\ \forall i,\ \forall \text{ primes } q<P\Big\}.\]
(The gcd condition is vacuous when \(|\mathcal F|=1\), so
\(S(g,P)\ge 1\).) Put \(F(g):=S(g,2)\) (no coprimality condition). The
statement of the problem is: \emph{every pairwise disjoint family of
size \(k\) has a pair with \(\gcd\ge k\)}; equivalently \(F(g)\le g\)
for all \(g\). The sharp example \(1,\dots,k \pmod k\) shows
\(F(g)\ge g\). We prove upper bounds for \(F\).

\subparagraph{Lemma 1 (two classes)}\label{lemma-1-two-classes}

If \(|\mathcal F|\ge 2\) and \(\mathcal F\) is pairwise disjoint then
\(\gcd(m_i,m_j)\ge 2\) for all \(i\ne j\).

\emph{Proof.} If \(\gcd(m_i,m_j)=1\) then \(1\mid a_i-a_j\), so by
\((*)\) the classes meet. \(\square\)

\subparagraph{Lemma 2 (base of the
recursion)}\label{lemma-2-base-of-the-recursion}

If \(P>g\) then \(S(g,P)=1\). In particular \(S(1,P)=1\) for every prime
\(P\).

\emph{Proof.} The one-element family \(\{(0,1)\}\) satisfies the
conditions, so \(S(g,P)\ge 1\). Suppose \(|\mathcal F|\ge 2\) satisfies
the conditions and take \(i\ne j\). By Lemma 1, \(d=\gcd(m_i,m_j)\ge 2\)
has a prime factor \(q\), with \(q\le d\le g<P\). But
\(q\mid d\mid m_i\) contradicts \(\gcd(m_i,q)=1\) for the prime \(q<P\).
Hence every admissible family has size \(1\). For \(g=1\) there is no
prime \(\le 1\), so \(P>g\) for all \(P\). \(\square\)

\subparagraph{\texorpdfstring{Lemma 3 (splitting inside one residue
class mod
\(p\))}{Lemma 3 (splitting inside one residue class mod p)}}\label{lemma-3-splitting-inside-one-residue-class-mod-p}

Let \(\mathcal F\) be pairwise disjoint, \(p\) a prime,
\(r\in\{0,\dots,p-1\}\), and
\[J_{p,r}=\{i\in I:\ p\mid m_i,\ a_i\equiv r \pmod p\}.\] For
\(i\in J_{p,r}\) put \(n_i=m_i/p\in\mathbb Z_{\ge1}\) and
\(b_i=(a_i-r)/p\in\mathbb Z\). Then: 1. the family
\(\mathcal F'=\{(b_i,n_i)\}_{i\in J_{p,r}}\) is pairwise disjoint; 2.
\(\gcd(n_i,n_j)=\gcd(m_i,m_j)/p\) for \(i\ne j\) in \(J_{p,r}\); 3. if
every \(m_i\) is coprime to all primes \(q<P\), so is every \(n_i\).

\emph{Proof.} Let \(\varphi:\mathbb Z\to\mathbb Z\),
\(\varphi(y)=r+py\); \(\varphi\) is injective. For \(i\in J_{p,r}\),
\[\varphi\big(b_i+n_i\mathbb Z\big)=r+p b_i+p n_i\mathbb Z=a_i+m_i\mathbb Z .\]
If \(y\in (b_i+n_i\mathbb Z)\cap(b_j+n_j\mathbb Z)\) then \(\varphi(y)\)
lies in both \(a_i+m_i\mathbb Z\) and \(a_j+m_j\mathbb Z\),
contradicting disjointness; this proves 1. For 2:
\(\gcd(m_i,m_j)=\gcd(pn_i,pn_j)=p\gcd(n_i,n_j)\). For 3:
\(n_i\mid m_i\), so any prime dividing \(n_i\) divides \(m_i\).
\(\square\)

\subparagraph{Lemma 4 (prime-splitting
recursion)}\label{lemma-4-prime-splitting-recursion}

Let \(g\ge 1\) and let \(P\) be a prime. Let \(p_1<p_2<\dots<p_s\) be
the primes in \([P,g]\) (\(s=0\) if \(P>g\)). Then
\[S(g,P)\ \le\ R(g,P):=\max\Big\{1,\ \sum_{t=1}^{s} p_t\, S\big(\lfloor g/p_t\rfloor,\,p_t\big)\Big\}.\]
Consequently \(S(g,P)\) is finite for all \(g,P\) (induction on \(g\),
since \(\lfloor g/p_t\rfloor<g\)).

\emph{Proof.} If \(s=0\) this is Lemma 2. Let \(s\ge1\) and let
\(\mathcal F\) be admissible for \(S(g,P)\) with \(|\mathcal F|\ge 2\)
(if all admissible families have size \(\le1\) the bound is trivial).
Define, for \(t=1,\dots,s\),
\[K_t=\{i\in I:\ p_t\mid m_i,\ p_u\nmid m_i \text{ for all } u<t\},\qquad K_{s+1}=\{i\in I: p_t\nmid m_i \text{ for all } t\le s\}.\]
These sets partition \(I\) (each \(i\) goes to the smallest \(t\) with
\(p_t\mid m_i\), or to \(K_{s+1}\)).

\emph{Claim: \(K_{s+1}=\emptyset\).} Let \(i\in K_{s+1}\) and pick
\(j\ne i\) (possible since \(|\mathcal F|\ge2\)). By Lemma 1,
\(\gcd(m_i,m_j)\ge2\) has a prime factor \(q\le\gcd(m_i,m_j)\le g\), and
\(q\mid m_i\). By admissibility \(q\ge P\) (no prime \(<P\) divides
\(m_i\)). So \(q\in\{p_1,\dots,p_s\}\) divides \(m_i\), contradicting
\(i\in K_{s+1}\).

\emph{Bound for \(K_t\).} Fix \(t\le s\) and write \(p=p_t\). Every
\(i\in K_t\) has \(p\mid m_i\), so
\(K_t=\bigsqcup_{r=0}^{p-1}(K_t\cap J_{p,r})\) with \(J_{p,r}\) as in
Lemma 3. Fix \(r\) and consider the sub-family
\(\mathcal F'=\{(b_i,n_i)\}_{i\in K_t\cap J_{p,r}}\) of Lemma 3. By
Lemma 3(1) it is pairwise disjoint. By Lemma 3(2), for \(i\ne j\) in it,
\(\gcd(n_i,n_j)=\gcd(m_i,m_j)/p\le g/p\), and being an integer,
\(\gcd(n_i,n_j)\le\lfloor g/p\rfloor\). Every \(m_i\) with \(i\in K_t\)
is coprime to \(p_1,\dots,p_{t-1}\) (definition of \(K_t\)) and to all
primes \(<P\) (admissibility), i.e.~to all primes \(<p_t\); by Lemma
3(3) the same holds for \(n_i\). Hence \(\mathcal F'\) is admissible for
\(S(\lfloor g/p_t\rfloor,p_t)\), and
\(|K_t\cap J_{p,r}|\le S(\lfloor g/p_t\rfloor,p_t)\). Summing over the
\(p_t\) values of \(r\): \(|K_t|\le p_t\,S(\lfloor g/p_t\rfloor,p_t)\).

Therefore
\(|\mathcal F|=\sum_{t=1}^{s}|K_t|\le\sum_{t=1}^{s}p_t\,S(\lfloor g/p_t\rfloor,p_t)\).
Taking the sup over \(\mathcal F\) gives the lemma. Finiteness: by
induction on \(g\); for \(g=1\) Lemma 2 gives \(1\); for \(g\ge2\) the
right side involves only \(S(g',\cdot)\) with
\(g'=\lfloor g/p_t\rfloor\le g/2<g\). \(\square\)

\subparagraph{Theorem 5 (closed bounds)}\label{theorem-5-closed-bounds}

For every \(g\ge1\) and every prime \(P\):
\[\text{(A)}\quad S(g,P)\le g^{3};\qquad\qquad \text{(B)}\quad S(g,P)\le g^{2}\prod_{\substack{p \text{ prime}\\ P\le p\le g}}\Big(1+\frac{2}{p}\Big).\]
In particular \(F(g)\le g^3\) and \(F(g)\le g^2\prod_{p\le g}(1+2/p)\).

\emph{Proof of (A).} Induction on \(g\). For \(g=1\): \(S(1,P)=1=1^3\)
(Lemma 2). Let \(g\ge2\) and assume (A) for all \(g'<g\) and all primes.
If no prime lies in \([P,g]\), Lemma 2 gives \(S=1\le g^3\). Otherwise,
by Lemma 4 and the induction hypothesis (note \(\lfloor g/p\rfloor<g\)),
\[S(g,P)\le\sum_{P\le p\le g}p\,\lfloor g/p\rfloor^{3}\le\sum_{P\le p\le g}p\Big(\frac gp\Big)^{3}=g^{3}\sum_{P\le p\le g}\frac1{p^{2}}\le g^{3}\sum_{n\ge2}\frac1{n^{2}}< g^{3}\sum_{n\ge2}\frac1{n(n-1)}=g^{3},\]
using \(\frac1{n^2}<\frac1{n(n-1)}\) and the telescoping sum
\(\sum_{n\ge2}\big(\frac1{n-1}-\frac1n\big)=1\). (Also \(R(g,P)\ge 1\)
trivially \(\le g^3\).) \(\square\)

\emph{Proof of (B).} Write
\(Q(g,P):=\prod_{P\le q\le g,\ q\text{ prime}}(1+2/q)\) (empty product
\(=1\)). Induction on \(g\). For \(g=1\): \(S(1,P)=1=1^2\cdot 1\). Let
\(g\ge2\), assume (B) for all \(g'<g\) and all primes. If there is no
prime in \([P,g]\), \(S(g,P)=1\le g^2\). Otherwise let \(p_1<\dots<p_s\)
be the primes in \([P,g]\) and set
\[Q_t:=\prod_{u=t}^{s}\Big(1+\frac2{p_u}\Big)\quad(1\le t\le s),\qquad Q_{s+1}:=1,\]
so \(Q_1=Q(g,P)\) and \(Q_t=(1+2/p_t)Q_{t+1}\). By Lemma 4 and the
induction hypothesis,
\[S(g,P)\le\sum_{t=1}^{s}p_t\,\lfloor g/p_t\rfloor^{2}\,Q\big(\lfloor g/p_t\rfloor,p_t\big).\]
Now \(\lfloor g/p_t\rfloor^2\le (g/p_t)^2\), and
\(Q(\lfloor g/p_t\rfloor,p_t)\le Q(g,p_t)=Q_t\) because the first
product runs over a subset of the primes of the second and every factor
is \(\ge1\). Hence \[S(g,P)\le g^{2}\sum_{t=1}^{s}\frac{Q_t}{p_t}.\] For
each \(t\), since \(p_t\ge2\) we have \(1+2/p_t\le 2\), so
\[\frac{Q_t}{p_t}=\frac{(1+2/p_t)\,Q_{t+1}}{p_t}\le\frac{2}{p_t}Q_{t+1}=Q_t-Q_{t+1}.\]
Summing, \(\sum_{t=1}^{s}Q_t/p_t\le Q_1-Q_{s+1}=Q_1-1<Q_1\). Therefore
\(S(g,P)\le g^2Q_1=g^2Q(g,P)\). \(\square\)

\subparagraph{Corollary 6 (lower bounds for the maximal
gcd)}\label{corollary-6-lower-bounds-for-the-maximal-gcd}

Let \(k\ge2\) and let \(a_1\pmod{m_1},\dots,a_k\pmod{m_k}\) be pairwise
disjoint, \(g:=\max_{i<j}\gcd(m_i,m_j)\). Then:

\begin{enumerate}
\def\labelenumi{\arabic{enumi}.}
\tightlist
\item
  (elementary, unconditional) \(g\ge k^{1/3}\);
\item
  (with Mertens' theorem, CITED) \(g\ge c\,\sqrt{k}/\log k\) for an
  absolute constant \(c>0\); explicitly, if one uses the
  Rosser--Schoenfeld inequality
  \(\sum_{p\le x}1/p<\log\log x+B+\frac{1}{2\log^2 x}\) (\(x>1\),
  \(B=0.2614\ldots\)), then
  \(g\ge\min\{\sqrt k,\ \sqrt k/(4\log k)\}=\sqrt k/(4\log k)\) for
  \(k\ge 2\) (see gap list: this constant depends on the cited
  inequality, which I have not re-derived).
\end{enumerate}

\emph{Proof.} The family is admissible for \(S(g,2)=F(g)\), so
\(k\le F(g)\). (1): \(k\le g^3\) by Theorem 5(A). (2): By Theorem 5(B),
\(k\le g^2\prod_{p\le g}(1+2/p)\le g^2\exp\big(2\sum_{p\le g}1/p\big)\),
using \(1+x\le e^x\). By Mertens' theorem
\(\sum_{p\le g}1/p=\log\log g+O(1)\), so \(k\le C\,g^2(\log g)^2\) for
an absolute \(C\) and all \(g\ge2\) (note \(g\ge2\) by Lemma 1). If
\(g\ge\sqrt k\) we are done. Otherwise \(g<\sqrt k\), so
\(\log g<\tfrac12\log k\) and \(k\le C g^2(\log k)^2/4\),
i.e.~\(g\ge (2/\sqrt C)\,\sqrt k/\log k\). With the Rosser--Schoenfeld
form, for \(g\ge2\):
\(\exp(2\sum_{p\le g}1/p)\le (\log g)^2 e^{2B+1/\log^2 g}\le (\log g)^2e^{2B+1/\log^2 2}<13.53(\log g)^2\),
so \(C=13.53\) and \(2/\sqrt{C}>0.544>1/4\). \(\square\)

\subparagraph{Remarks (what the method gives and where it
loses)}\label{remarks-what-the-method-gives-and-where-it-loses}

\begin{itemize}
\tightlist
\item
  \(R(2,2)=2\): every pairwise disjoint family with all gcds \(\le2\)
  has at most \(2\) members, i.e.~\textbf{any three pairwise disjoint
  classes have a pair with \(\gcd\ge3\)} --- the statement of cell 1
  falls out of Lemma 4 alone. \(R(3,2)=5>3\), so cell 2 does not.
\item
  The exact values of the recursion \(R(g,2)\) (integer arithmetic,
  script below, \(g\le3000\), 2.3 s) satisfy \(R(g,2)\le g^2\) for all
  \(1\le g\le 3000\), with \(R(g,2)/g^2\to\) about \(0.50\)
  (e.g.~\(R(1000,2)=494228\)). So the recursion itself cannot give
  better than \(F(g)\lesssim g^2/2\): the loss relative to the
  conjecture \(F(g)=g\) (and to the \(k^{1-o(1)}\) of arXiv:2607.24655)
  sits entirely in the union bound \(|\mathcal F|\le\sum_t|K_t|\) of
  Lemma 4, which ignores that the ``\(p\)-cliques'' \(\{i:p\mid m_i\}\)
  for different primes \(p\) must overlap heavily (every pair of
  vertices shares a prime). Exploiting this overlap is the natural next
  subgoal; the closed bounds (A),(B) are asymptotically far from
  \(R(g,2)\) (Theorem 5(B) gives \(\approx 4.4\cdot10^7\) at \(g=1000\)
  vs \(R=4.9\cdot10^5\)), so there is also room to sharpen the induction
  to \(F(g)\le g^2\) without touching Lemma 4.
\item
  Everything above uses only \((*)\); no result from the literature
  enters Lemmas 1--4 or Theorem 5. Mertens/Rosser--Schoenfeld enters
  only Corollary 6(2).
\end{itemize}


% ===== CAMPO: 3. Verification: instructions, dependencies, timings =====
Python 3 standard library. Scripts (re-run by the orchestrator, see
\texttt{runs/p4\_c6/verifica/}): - code\_1 (python, rigor
\texttt{exact}): Exact integer evaluation of the recursion R(g,P) of
Lemma 4 for all g \textless= 3000 (P ranging over primes), and
exact-rational check that R(g,2) \textless= g\^{}2 *
prod\_\{p\textless=g\}(1+2/p) and R(g,2) \textless= g\^{}3 hold for
every g \textless= 3000 (consistency check of Theorem 5 against the
recursion; not needed for the proof). Finite set covered: g in
{[}1,3000{]}. Wall clock: 2.34 s (Python 3, .venv). Command:
.venv/bin/python runs/p4\_c6/sandbox/ricorsione\_gcd.py 3000


% ===== CAMPO: 4. Sources and contribution =====
\begin{itemize}
\tightlist
\item
  arXiv:2607.24655v1, J. Fornal, Y.-C. Sun, `On the problem of large gcd
  for disjoint residue classes' (2026), Theorem 1.1: max gcd(m\_i,m\_j)
  \textgreater\textgreater{} k exp(-(2+o(1)) sqrt(log k/log log k)) ---
  CITED for context only (fetched; introduction and section structure
  read, proofs not reproduced)
\item
  K. O'Bryant, `On Z.-W. Sun's disjoint congruence classes conjecture',
  Combinatorial Number Theory, de Gruyter 2007, 403--412 --- CITED at
  second hand via the introduction of arXiv:2607.24655 (integer
  conjecture for k \textless= 20); not read
\item
  Z.-W. Sun, `Finite covers of groups by cosets or subgroups', Internat.
  J. Math. 17 (2006) 1047--1064 --- origin of the conjecture and its
  group form; cited at second hand via arXiv:2607.24655, not read
\item
  Mertens' theorem, sum\_\{p\textless=x\} 1/p = log log x + O(1);
  explicit form J. B. Rosser, L. Schoenfeld, Illinois J. Math. 6 (1962),
  Thm 5 (3.20): sum\_\{p\textless=x\}1/p \textless{} log log x + B +
  1/(2 log\^{}2 x), x\textgreater1, B = 0.2614972\ldots{} --- CITED from
  memory, used only in Corollary 6(2)
\item
  Position with respect to the literature: State of the art for this
  cell (direction (b)): arXiv:2607.24655 (Fornal--Sun, 27 Jul 2026),
  Theorem 1.1, which I fetched and read in part (introduction and
  section structure, not the full proofs): ``Let a\_1 (mod
  m\_1),\ldots,a\_k (mod m\_k) be pairwise disjoint residue classes.
  Then max gcd(m\_i,m\_j) \textgreater\textgreater{} k exp(-(2+o(1))
  sqrt(log k / log log k)).'' --- CITED, not reproduced. Note this is
  \emph{stronger} than the bound quoted in the problem statement
  (exponent log k/log log k without the square root): the square-root
  version already ``improves the exponential factor'' asked for in (b);
  but by the hand-in rules a citation scores nothing, and the paper is
  \textasciitilde2000 lines (weighted gcd graph, sieve-theoretic
  partition Prop. 2.1, Möbius inversion and DFT in Prop. 2.2, key Lemma
  4.1), which I judged not reproducible and checkable within one
  iteration. Its introduction also records (CITED, not checked):
  O'Bryant proved the integer conjecture for k \textless= 20 {[}K.
  O'Bryant, On Z.-W. Sun's disjoint congruence classes conjecture,
  Combinatorial Number Theory, de Gruyter 2007, 403--412{]}; Zhu proved
  the group form for k=3,4 {[}Int. J. Mod. Math. 3 (2008){]}; Sun proved
  k=2 and finite p-groups {[}Internat. J. Math. 17 (2006){]}. The
  problem's statement that the group form is known for k\textless=5 is
  consistent with arXiv:2608.15873 (Menon) / Margolis--Schnabel on
  G-harmonic tuples, relevant only to direction (c). arXiv:2603.26043
  and arXiv:1511.04293 concern disjoint \emph{covering} systems with one
  repeated modulus and are not relevant here. My approach departs from
  Fornal--Sun: instead of weighting vertices over all n \textless= d and
  using Fourier analysis, it uses only the affine rescaling y
  -\textgreater{} r+py inside a residue class mod a prime, an elementary
  recursion, and a closed-form induction. The price is the exponent (1/2
  instead of 1-o(1)); the gain is a complete proof.
\end{itemize}


% ===== CAMPO: 5. Limits and unresolved parts =====
Referee's blocking point: The strongest result actually proved is
Corollary 6: max gcd(m\_i,m\_j) \textgreater= k\^{}\{1/3\}
unconditionally, and \textgreater= c\emph{sqrt(k)/log k conditional on
Mertens/Rosser--Schoenfeld. The cell-6 target requires one of: (a)
deciding a size k \textgreater= 25; (b) gcd \textgreater= ck for an
absolute c \textgreater{} 0, or an improvement of the known factor
k}exp(-(2+o(1)) log k/log log k) (i.e.~something better than
k\^{}\{1-o(1)\}); (c) the group form at k = 6. A bound of order
k\^{}\{1/2-o(1)\} is strictly weaker than the already-known
k\^{}\{1-o(1)\} and does not touch (a) or (c). The submission's own
remarks concede that the recursion of Lemma 4 cannot give better than
F(g) ≲ g\^{}2/2, i.e.~gcd ≳ sqrt(2k), so no route to the target is
present. The declared claim `main' is therefore not established, and no
other reusable claim was declared. Next step required: Address the
stated blocking obligation without silently changing the target Gaps
declared by the author: - Corollary 6(2) relies on Mertens' theorem
(and, for the explicit constant 1/4, on the Rosser--Schoenfeld
inequality sum\_\{p\textless=x\}1/p \textless{} log log x + B + 1/(2
log\^{}2 x)), quoted from memory and not re-derived here; the
qualitative form `c sqrt(k)/log k for some absolute c' needs only
Mertens, the explicit constant needs the RS form. Everything else
(Lemmas 1--4, Theorem 5, Corollary 6(1)) is self-contained. - The result
is weaker than the literature (exponent 1/2 vs 1-o(1)); it does not
achieve any of the three targets of cell 6. The numerical observation
R(g,2) \textless= g\^{}2 for g \textless= 3000 is evidence only and is
not used in any proof. - Literature statements attributed to O'Bryant (k
\textless= 20), Zhu (group form k=3,4) and Sun (k=2, p-groups) are taken
from the introduction of arXiv:2607.24655 and were not checked against
the original papers.

